%% Abstract (IEEE Access: 150-250 words, one paragraph, no equations/citations) and index terms
\begin{abstract}
Transformerless single-phase inverters raise efficiency in grid-connected photovoltaic systems, but the missing galvanic isolation lets the switching common-mode voltage drive a leakage current through the parasitic capacitance between the photovoltaic array and ground. Safety standards require this current to stay below 300 mA. Most topology comparisons use a single switching frequency. This paper instead studies how the switching frequency interacts with the common-mode resonance formed by the output-filter inductance and the parasitic capacitance. Simulation models of the H4 inverter and the highly efficient and reliable inverter concept (HERIC), each with unipolar and bipolar sinusoidal pulse-width modulation, an inductive filter and a 150 nF parasitic capacitance, are swept from 1 kHz to 40 kHz while leakage current, efficiency and low-order harmonic distortion are recorded. In all four cases the leakage current peaks sharply, reaching about 0.75 A, 0.5 A and 4 A for H4 unipolar, H4 bipolar and HERIC unipolar, at frequencies that, for the dominant peaks, agree within about 5\% with the analytically predicted resonance, its sub-multiples or, in one case, twice its value. Elsewhere the leakage falls to 3.5--40 mA. With identical circuit parameters, bipolar HERIC leaks about eleven times less than unipolar HERIC, but its low-order distortion of about 9.1\% is high. The results are simulation-based and the H4 and HERIC parameter sets differ, so they support a switching-frequency selection rule (avoid the resonance, its sub-multiples and twice its value; prefer frequencies above 2.4 times the resonance) rather than a topology ranking.
\end{abstract}

\begin{keywords}
Common-mode voltage, H4 inverter, HERIC inverter, leakage current, parasitic capacitance, photovoltaic systems, switching frequency, transformerless inverter.
\end{keywords}
